\documentclass[10pt]{extarticle}
\usepackage[a4paper,margin=1.5cm]{geometry}
\usepackage[T1]{fontenc}
\usepackage{amsmath,amssymb}
\usepackage{graphicx}
\usepackage[export]{adjustbox}
\usepackage{multicol}
\usepackage{enumitem}
\usepackage{booktabs,array,tabularx}
\setlength{\columnsep}{0.7cm}
\setlength{\parindent}{0pt}
\setlength{\parskip}{3pt}
\pagestyle{empty}
\makeatletter
% figure registry (figdefs.tex is generated by the extract step)
\newcommand{\qfigdef}[4]{\@namedef{qf@file@#1}{#2}\@namedef{qf@w@#1}{#3}\@namedef{qf@wide@#1}{#4}}
\newcommand{\qf@check}[1]{\@ifundefined{qf@file@#1}{\PackageError{score2}{Unknown figure id #1}{}}{}}
\newcommand{\qf@img}[4][]{\includegraphics[width=#2\dimexpr\@nameuse{qf@w@#4}pt\relax,max width=#3,max height=0.45\textheight,#1]{\@nameuse{qf@file@#4}}}
% Full-width questions: finalize wraps every question that contains a wide figure in \qfullbegin...\qfullend,
% so the WHOLE question (stem, figures, options) is set across the page in reading order, then the two columns
% resume. The question's head - its number, stem and first figure (\qfullsplit marks the end of the head) -
% is kept on one page, so a number is never left at a page foot with its figure overleaf; the rest may break
% across pages normally, so no big gaps are left. Without a \qfullsplit, the whole question is kept together
% if it is under 45% of a page. Never write these yourself.
\newif\ifqf@full
\newif\ifqf@boxed
\newsavebox{\qf@box}
\newcommand{\qf@open}{\global\qf@fulltrue\par\global\qf@boxedtrue
  \setbox\qf@box\vbox\bgroup\hsize\textwidth\linewidth\textwidth}
\newcommand{\qf@emit}[1]{\par\egroup\global\qf@boxedfalse
  \ifdim\dimexpr\ht\qf@box+\dp\qf@box\relax<#1\textheight\noindent\usebox{\qf@box}\par\else\unvbox\qf@box\fi}
\newcommand{\qfullsplit}{\ifqf@boxed\qf@emit{0.85}\fi}
\newcommand{\qf@close}{\ifqf@boxed\qf@emit{0.45}\fi\global\qf@fullfalse}
\newcommand{\qfullbegin}{\end{multicols}\qf@open}
\newcommand{\qfullend}{\qf@close\begin{multicols}{2}}
\newcommand{\qfullstart}{\qf@open}
\newcommand{\qfulldone}{\qf@close}
% \qfig[scale]{id}: block figure, centred, directly where it is written (size set by the tool).
% Fallback if a wide figure is not inside a full-width question: break out of the columns in place - no floats.
\newcommand{\qfig}[2][1]{\qf@check{#2}%
  \ifqf@full\qf@blk{#1}{#2}\else
  \ifnum\@nameuse{qf@wide@#2}=1 \expandafter\@firstoftwo\else\expandafter\@secondoftwo\fi
  {\end{multicols}\begin{center}\qf@img{#1}{\textwidth}{#2}\end{center}\begin{multicols}{2}}%
  {\qf@blk{#1}{#2}}\fi}
\newcommand{\qf@blk}[2]{\par\smallskip{\centering\qf@img{#1}{\linewidth}{#2}\par}\smallskip}
% \qfigi[scale]{id}: inline figure (inside option lists / tables / minipages); never breaks columns.
\newcommand{\qfigi}[2][1]{\qf@check{#2}\qf@img[valign=c]{#1}{\linewidth}{#2}}
\makeatother
% structure macros - lint reads these to check completeness and figure placement
\newcommand{\papersec}[1]{\par\bigskip{\centering\Large\bfseries PAPER-#1\par}\medskip}
\newcommand{\partsec}[2]{\par\medskip\noindent{\large\bfseries PART-#1: #2}\par\smallskip}
\newcommand{\sect}[1]{\par\smallskip\noindent\textbf{#1}\par}
\newcommand{\qn}[1]{\par\smallskip\noindent\textbf{#1.}~}
\newcommand{\paraq}[2]{\par\smallskip\noindent\textbf{Paragraph for Questions #1 and #2}\par}
\newenvironment{qoptions}{\begin{itemize}[leftmargin=1.8em,itemsep=1pt,topsep=2pt,label={}]}{\end{itemize}}
\newcommand{\opt}[2]{\item[(#1)] #2}
% figure registry (generated by score2_tools.py finalize)
\qfigdef{p1_phy_s1_q01_1}{TOAS - Test 17 (18 Sep 2022)-img/p1_phy_s1_q01_1.png}{161.1}{0}
\qfigdef{p1_phy_s1_q02_1}{TOAS - Test 17 (18 Sep 2022)-img/p1_phy_s1_q02_1.png}{134.5}{0}
\qfigdef{p1_phy_s1_q03_1}{TOAS - Test 17 (18 Sep 2022)-img/p1_phy_s1_q03_1.png}{161.1}{0}
\qfigdef{p1_phy_s1_q05_1}{TOAS - Test 17 (18 Sep 2022)-img/p1_phy_s1_q05_1.png}{266.7}{1}
\qfigdef{p1_phy_s1_q07_1}{TOAS - Test 17 (18 Sep 2022)-img/p1_phy_s1_q07_1.png}{369.7}{1}
\qfigdef{p1_phy_s2_para11_13_1}{TOAS - Test 17 (18 Sep 2022)-img/p1_phy_s2_para11_13_1.png}{55.6}{0}
\qfigdef{p1_phy_s2_para11_13_5}{TOAS - Test 17 (18 Sep 2022)-img/p1_phy_s2_para11_13_5.png}{35.1}{0}
\qfigdef{p1_phy_s2_para11_13_2}{TOAS - Test 17 (18 Sep 2022)-img/p1_phy_s2_para11_13_2.png}{55.6}{0}
\qfigdef{p1_phy_s2_para11_13_6}{TOAS - Test 17 (18 Sep 2022)-img/p1_phy_s2_para11_13_6.png}{60.5}{0}
\qfigdef{p1_phy_s2_para11_13_3}{TOAS - Test 17 (18 Sep 2022)-img/p1_phy_s2_para11_13_3.png}{55.6}{0}
\qfigdef{p1_phy_s2_para11_13_7}{TOAS - Test 17 (18 Sep 2022)-img/p1_phy_s2_para11_13_7.png}{61.2}{0}
\qfigdef{p1_phy_s2_para11_13_4}{TOAS - Test 17 (18 Sep 2022)-img/p1_phy_s2_para11_13_4.png}{51.4}{0}
\qfigdef{p1_phy_s2_para11_13_8}{TOAS - Test 17 (18 Sep 2022)-img/p1_phy_s2_para11_13_8.png}{52.9}{0}
\qfigdef{p1_phy_s3_q01_1}{TOAS - Test 17 (18 Sep 2022)-img/p1_phy_s3_q01_1.png}{134.5}{0}
\qfigdef{p1_phy_s3_q05_1}{TOAS - Test 17 (18 Sep 2022)-img/p1_phy_s3_q05_1.png}{108.3}{0}
\qfigdef{p1_phy_s3_q05_2}{TOAS - Test 17 (18 Sep 2022)-img/p1_phy_s3_q05_2.png}{130.5}{0}
\qfigdef{p1_chem_s1_q01_1}{TOAS - Test 17 (18 Sep 2022)-img/p1_chem_s1_q01_1.png}{245.2}{0}
\qfigdef{p1_chem_s1_q02_1}{TOAS - Test 17 (18 Sep 2022)-img/p1_chem_s1_q02_1.png}{210.6}{0}
\qfigdef{p1_chem_s1_q03_1}{TOAS - Test 17 (18 Sep 2022)-img/p1_chem_s1_q03_1.png}{174.8}{0}
\qfigdef{p1_chem_s2_pre_1}{TOAS - Test 17 (18 Sep 2022)-img/p1_chem_s2_pre_1.png}{45.1}{0}
\qfigdef{p1_chem_s2_pre_2}{TOAS - Test 17 (18 Sep 2022)-img/p1_chem_s2_pre_2.png}{62.4}{0}
\qfigdef{p1_chem_s2_pre_3}{TOAS - Test 17 (18 Sep 2022)-img/p1_chem_s2_pre_3.png}{69.2}{0}
\qfigdef{p1_chem_s2_pre_4}{TOAS - Test 17 (18 Sep 2022)-img/p1_chem_s2_pre_4.png}{69.2}{0}
\qfigdef{p1_chem_s3_q03_1}{TOAS - Test 17 (18 Sep 2022)-img/p1_chem_s3_q03_1.png}{204.5}{0}
\qfigdef{p2_phy_s1_q09_1}{TOAS - Test 17 (18 Sep 2022)-img/p2_phy_s1_q09_1.png}{148.7}{0}
\qfigdef{p2_phy_s1_q12_1}{TOAS - Test 17 (18 Sep 2022)-img/p2_phy_s1_q12_1.png}{100.4}{0}
\qfigdef{p2_phy_s1_q13_1}{TOAS - Test 17 (18 Sep 2022)-img/p2_phy_s1_q13_1.png}{42.4}{0}
\qfigdef{p2_phy_s1_q13_A}{TOAS - Test 17 (18 Sep 2022)-img/p2_phy_s1_q13_A.png}{81.1}{0}
\qfigdef{p2_phy_s1_q13_B}{TOAS - Test 17 (18 Sep 2022)-img/p2_phy_s1_q13_B.png}{81.1}{0}
\qfigdef{p2_phy_s1_q13_C}{TOAS - Test 17 (18 Sep 2022)-img/p2_phy_s1_q13_C.png}{89.1}{0}
\qfigdef{p2_phy_s1_q14_1}{TOAS - Test 17 (18 Sep 2022)-img/p2_phy_s1_q14_1.png}{154.5}{0}
\qfigdef{p2_phy_s1_q15_1}{TOAS - Test 17 (18 Sep 2022)-img/p2_phy_s1_q15_1.png}{133.7}{0}
\qfigdef{p2_phy_s1_q16_1}{TOAS - Test 17 (18 Sep 2022)-img/p2_phy_s1_q16_1.png}{108.3}{0}
\qfigdef{p2_phy_s1_q17_A}{TOAS - Test 17 (18 Sep 2022)-img/p2_phy_s1_q17_A.png}{62.0}{0}
\qfigdef{p2_phy_s1_q17_B}{TOAS - Test 17 (18 Sep 2022)-img/p2_phy_s1_q17_B.png}{62.0}{0}
\qfigdef{p2_phy_s1_q17_C}{TOAS - Test 17 (18 Sep 2022)-img/p2_phy_s1_q17_C.png}{61.8}{0}
\qfigdef{p2_phy_s1_q17_D}{TOAS - Test 17 (18 Sep 2022)-img/p2_phy_s1_q17_D.png}{61.8}{0}
\qfigdef{p2_phy_s1_q19_A}{TOAS - Test 17 (18 Sep 2022)-img/p2_phy_s1_q19_A.png}{74.0}{0}
\qfigdef{p2_phy_s1_q19_B}{TOAS - Test 17 (18 Sep 2022)-img/p2_phy_s1_q19_B.png}{74.0}{0}
\qfigdef{p2_phy_s1_q19_C}{TOAS - Test 17 (18 Sep 2022)-img/p2_phy_s1_q19_C.png}{74.0}{0}
\qfigdef{p2_phy_s1_q19_D}{TOAS - Test 17 (18 Sep 2022)-img/p2_phy_s1_q19_D.png}{74.0}{0}
\qfigdef{p2_phy_s2_q01_1}{TOAS - Test 17 (18 Sep 2022)-img/p2_phy_s2_q01_1.png}{139.9}{0}
\qfigdef{p2_phy_s2_q02_1}{TOAS - Test 17 (18 Sep 2022)-img/p2_phy_s2_q02_1.png}{128.3}{0}
\qfigdef{p2_phy_s2_q04_1}{TOAS - Test 17 (18 Sep 2022)-img/p2_phy_s2_q04_1.png}{112.4}{0}
\qfigdef{p2_phy_s2_q07_1}{TOAS - Test 17 (18 Sep 2022)-img/p2_phy_s2_q07_1.png}{107.8}{0}
\qfigdef{p2_phy_s2_q08_1}{TOAS - Test 17 (18 Sep 2022)-img/p2_phy_s2_q08_1.png}{154.2}{0}
\qfigdef{p2_chem_s1_q05_1}{TOAS - Test 17 (18 Sep 2022)-img/p2_chem_s1_q05_1.png}{112.3}{0}
\qfigdef{p2_chem_s1_q06_1}{TOAS - Test 17 (18 Sep 2022)-img/p2_chem_s1_q06_1.png}{139.6}{0}
\qfigdef{p2_chem_s1_q06_A}{TOAS - Test 17 (18 Sep 2022)-img/p2_chem_s1_q06_A.png}{55.6}{0}
\qfigdef{p2_chem_s1_q06_B}{TOAS - Test 17 (18 Sep 2022)-img/p2_chem_s1_q06_B.png}{63.0}{0}
\qfigdef{p2_chem_s1_q06_C}{TOAS - Test 17 (18 Sep 2022)-img/p2_chem_s1_q06_C.png}{56.5}{0}
\qfigdef{p2_chem_s1_q06_D}{TOAS - Test 17 (18 Sep 2022)-img/p2_chem_s1_q06_D.png}{61.0}{0}
\qfigdef{p2_chem_s1_q07_D}{TOAS - Test 17 (18 Sep 2022)-img/p2_chem_s1_q07_D.png}{45.1}{0}
\qfigdef{p2_chem_s1_q08_2}{TOAS - Test 17 (18 Sep 2022)-img/p2_chem_s1_q08_2.png}{24.0}{0}
\qfigdef{p2_chem_s1_q08_C}{TOAS - Test 17 (18 Sep 2022)-img/p2_chem_s1_q08_C.png}{128.3}{0}
\qfigdef{p2_chem_s1_q08_D}{TOAS - Test 17 (18 Sep 2022)-img/p2_chem_s1_q08_D.png}{83.5}{0}
\qfigdef{p2_chem_s1_q09_A}{TOAS - Test 17 (18 Sep 2022)-img/p2_chem_s1_q09_A.png}{41.9}{0}
\qfigdef{p2_chem_s1_q09_B}{TOAS - Test 17 (18 Sep 2022)-img/p2_chem_s1_q09_B.png}{43.6}{0}
\qfigdef{p2_chem_s1_q09_C}{TOAS - Test 17 (18 Sep 2022)-img/p2_chem_s1_q09_C.png}{47.5}{0}
\qfigdef{p2_chem_s1_q09_D}{TOAS - Test 17 (18 Sep 2022)-img/p2_chem_s1_q09_D.png}{42.6}{0}
\qfigdef{p2_chem_s1_q10_B}{TOAS - Test 17 (18 Sep 2022)-img/p2_chem_s1_q10_B.png}{68.9}{0}
\qfigdef{p2_chem_s1_q11_A}{TOAS - Test 17 (18 Sep 2022)-img/p2_chem_s1_q11_A.png}{91.3}{0}
\qfigdef{p2_chem_s1_q11_B}{TOAS - Test 17 (18 Sep 2022)-img/p2_chem_s1_q11_B.png}{74.5}{0}
\qfigdef{p2_chem_s1_q11_C}{TOAS - Test 17 (18 Sep 2022)-img/p2_chem_s1_q11_C.png}{94.8}{0}
\qfigdef{p2_chem_s1_q11_D}{TOAS - Test 17 (18 Sep 2022)-img/p2_chem_s1_q11_D.png}{75.7}{0}
\qfigdef{p2_chem_s1_q12_1}{TOAS - Test 17 (18 Sep 2022)-img/p2_chem_s1_q12_1.png}{89.2}{0}
\qfigdef{p2_chem_s1_q12_A}{TOAS - Test 17 (18 Sep 2022)-img/p2_chem_s1_q12_A.png}{95.8}{0}
\qfigdef{p2_chem_s1_q12_B}{TOAS - Test 17 (18 Sep 2022)-img/p2_chem_s1_q12_B.png}{91.8}{0}
\qfigdef{p2_chem_s1_q12_C}{TOAS - Test 17 (18 Sep 2022)-img/p2_chem_s1_q12_C.png}{72.2}{0}
\qfigdef{p2_chem_s1_q12_D}{TOAS - Test 17 (18 Sep 2022)-img/p2_chem_s1_q12_D.png}{86.9}{0}
\qfigdef{p2_chem_s1_q13_1}{TOAS - Test 17 (18 Sep 2022)-img/p2_chem_s1_q13_1.png}{113.2}{0}
\qfigdef{p2_chem_s1_q13_B}{TOAS - Test 17 (18 Sep 2022)-img/p2_chem_s1_q13_B.png}{138.7}{0}
\qfigdef{p2_chem_s1_q13_C}{TOAS - Test 17 (18 Sep 2022)-img/p2_chem_s1_q13_C.png}{138.7}{0}
\qfigdef{p2_chem_s1_q13_D}{TOAS - Test 17 (18 Sep 2022)-img/p2_chem_s1_q13_D.png}{134.5}{0}
\qfigdef{p2_chem_s2_q01_1}{TOAS - Test 17 (18 Sep 2022)-img/p2_chem_s2_q01_1.png}{152.0}{0}
\qfigdef{p2_chem_s2_q02_1}{TOAS - Test 17 (18 Sep 2022)-img/p2_chem_s2_q02_1.png}{130.1}{0}
\qfigdef{p2_chem_s2_q06_2}{TOAS - Test 17 (18 Sep 2022)-img/p2_chem_s2_q06_2.png}{43.6}{0}
\qfigdef{p2_chem_s2_q08_1}{TOAS - Test 17 (18 Sep 2022)-img/p2_chem_s2_q08_1.png}{143.2}{0}
\raggedcolumns
\makeatletter\@beginparpenalty=10000\makeatother
\newcolumntype{L}[1]{>{\raggedright\arraybackslash}p{#1}}
\begin{document}
\begin{center}{\Large\bfseries TOAS --- Test 17 (18 Sep 2022)}\end{center}
\begin{multicols}{2}
\raggedcolumns
\papersec{1}
\partsec{1}{PHYSICS}
\sect{SECTION-I(i)}
\qn{1} A block of mass $M$ is connected to a uniform rope of length $L$ having uniform mass per unit length $\mu$. The other end of rope is hinged at $O$ and system is rotating with constant angular velocity $\omega$ in gravity free space. A transverse pulse of wavelength $\lambda$ is produced at $O$. Assume no loss in energy. Then (Given : $\mu=\dfrac{2M}{3L}$)
\qfig{p1_phy_s1_q01_1}
\begin{qoptions}
\opt{A}{Wavelength of pulse decrease as it reaches point $A$ and becomes $\dfrac{\sqrt3}{2}\lambda$.}
\opt{B}{Time taken by pulse to travel from $O$ to $A$ is $\dfrac{\pi}{3\sqrt2\,\omega}$.}
\opt{C}{As pulse propagate from $O$ to $A$, it's Amplitude continuously increases.}
\opt{D}{Pulse speed at point $A$ in frame of ground is $\sqrt{\dfrac52}\,\omega L$.}
\end{qoptions}
\qn{2} Four plane mirrors form a square with each side of length $L$. There is a small hole at one of the corners as shown in the figure. A ray of light passed through the hole at an angle $\theta$ and can come out of same hole. The cross-section of mirror configuration and the ray of light lies on the same plane. Which of the following statements are correct.
\qfig{p1_phy_s1_q02_1}
\begin{qoptions}
\opt{A}{A ray of light will come out for $\theta=45^\circ$ only if $\ell=\dfrac L2$.}
\opt{B}{A ray of light will come out for $\theta=90^\circ$ for $0<\ell<L$.}
\opt{C}{A ray of light will come out for $\theta=\tan^{-1}\left(\dfrac1{2n}\right)$, where $n$ is a natural number for $\ell=\dfrac L2$.}
\opt{D}{A ray of light will come out for $\theta=45^\circ$ for $0<\ell<L$.}
\end{qoptions}
\qn{3} In given setup coherent light of wave length $\lambda$ is used assuming $D\gg d$, select correct alternative [each slit produce same intensity on screen]
\qfig{p1_phy_s1_q03_1}
\begin{qoptions}
\opt{A}{For $d=\sqrt{2\lambda D}$, intensity at C and A will be maximum but minimum at B.}
\opt{B}{For $d=2\sqrt{\lambda D}$, intensity at A, B and C will be maximum intensity.}
\opt{C}{For $d=\sqrt{\dfrac{\lambda D}2}$, intensity at A, B and C will be $\left(\dfrac59\right)^{\text{th}}$ maximum intensity.}
\opt{D}{For $d=\dfrac{\sqrt{\lambda D}}2$, intensity at A and C will be $\left(\dfrac19\right)^{\text{th}}$ maximum intensity.}
\end{qoptions}
\qn{4} Two long parallel straight wire cavey equal current $I$ along positive $z$-direction. One of then intersect $x$-axis at point $(a,0,0)$ and other at $(-a,0,0)$. A third finite length wire of uniform mass per unit length $\lambda$ carries current in negative $z$-direction and intersect $y$-axis at point $(0,b,0)$. If acceleration due to gravity is $\vec g=-g\hat j$ and wire is suspended under gravitational and magnetic forces only then select correct alternatives.
\begin{qoptions}
\opt{A}{Current in finite length wire is $\dfrac{\pi\lambda g(a^2+b^2)}{\mu_0Ib}$.}
\opt{B}{When wire is displaced by $\Delta y$ along $y$-axis and released then wire will perform SHM for $\Delta y\ll b$.}
\opt{C}{Wire is displaced by $\Delta y$ along $y$-direction and released then wire will perform SHM for $\Delta y\ll b$ and $b>a$.}
\opt{D}{Wire will perform SHM if displaced slightly in any direction perpendicular to $z$-axis.}
\end{qoptions}
\qfullbegin
\qn{5} A rod of length $L=1$m is made of material having density $\rho=\dfrac{3g}{cc}$ and youngs modulus $Y=2.43\times10^9$ N/m$^2$, is clamped at distance $\dfrac L4$ from one end and other end has as metal disk attached. Glass tubs contain gas of density $\rho_0=0.01$ g/cc and bulk modulus $B_0=9\times10^5$ N/m$^2$. When longitudanal oscillation are setup in rod then fine powder present in tube pill up at regular interval of $\ell=3.33$ cm, then select correct alternatives.
\qfig{p1_phy_s1_q05_1}
\qfullsplit
\begin{qoptions}
\opt{A}{Speed of wave in rod is 900m/s.}
\opt{B}{Speed of sound in gas 300 m/s}
\opt{C}{frequency of wave is nearly 4500 Hz.}
\opt{D}{Rod is vibrating in 2$^{\text{nd}}$ overtone.}
\end{qoptions}
\qfullend
\qn{6} Consider an ideal solenoid having radius of cross section $R$. There are $n$ turns/length. It carries a current of $I_0\sin(\omega t)$. Consider a cylinder which is coaxial with the solenoid and has a length $\ell$ and radius $r$ we want to find the rate of energy flow through this cylinder. ($r<R$) neglect magnetic field of induced electric field
\begin{qoptions}
\opt{A}{The rate of energy flow through its flat face is zero.}
\opt{B}{At $t=\dfrac\pi{4\omega}$, energy is flowing outwards from the curved surface of the cylinder.}
\opt{C}{The rate of energy flow is proportional to $\mu_0n^2i_0^2r^2\omega\ell$.}
\opt{D}{At $t=\dfrac\pi{4\omega}$, rate of energy flow is maximum.}
\end{qoptions}
\qfullbegin
\qn{7} Consider a diffraction pattern of 2 slits as shown in figure. The pattern is formed by 2 slits of width $a=3\mu$m each separated by $d=18\ \mu$m.
\qfig{p1_phy_s1_q07_1}
\qfullsplit
\begin{qoptions}
\opt{A}{Here 6$^{\text{th}}$ order maxima of interference can't be seen.}
\opt{B}{If the wavelength is decreased, the central envelope becomes narrower.}
\opt{C}{If the wavelength is decreased, the width of each interference maxima becomes lesser.}
\opt{D}{If the wavelength is decreased, same number of interference maxima will be seen in the central envelope of diffraction.}
\end{qoptions}
\qfullend
\sect{SECTION-I(ii)}
\paraq{8}{10}
A charged particle in introduced at the origin ($x=0$, $y=0$, $z=0$) with a given initial velocity $\vec V$. A uniform electric field $\vec E$ exist every where. The velocity $\vec V$, electric field $\vec E$ and magnetic field $\vec B$ are given in columns 1, 2 and 3 respectively. The quantities $E_0$, $B_0$ are positive in magnitude.\par
\begin{center}\footnotesize\setlength{\tabcolsep}{2pt}
\begin{tabular}{|L{3.4cm}|L{1.5cm}|L{1.5cm}|}
\hline
\multicolumn{1}{|c|}{Column-1}&\multicolumn{1}{c|}{Column-2}&\multicolumn{1}{c|}{Column-3}\\\hline
(I) Positive charge with $\vec V=\dfrac{2E_0}{B_0}\hat i$.&(i) $\vec E=E_0\hat j$&(P) $\vec B=-B_0\hat k$\\\hline
(II) Negative charge with $\vec V=\dfrac{E_0}{B_0}(\hat j+\hat k)$.&(ii) $\vec E=E_0\hat i$&(Q) $\vec B=B_0\hat i$\\\hline
(III) Positive charge with $\vec V=\dfrac{E_0}{B_0}(\hat i+\hat j)$.&(iii) $\vec E=E_0\hat k$&(R) $\vec B=-B_0\hat j$\\\hline
(IV) Negative charge with $\vec V=\dfrac{E_0}{B_0}\hat i$.&(iv) $\vec E=-E_0\hat i$&(S) $\vec B=B_0\hat k$\\\hline
\end{tabular}\end{center}
\qn{8} In which case particle will move on a helical path with decreasing pitch.
\begin{qoptions}
\opt{A}{(IV), (ii), (Q)}
\opt{B}{(I), (iii), (R)}
\opt{C}{(II), (i), (R)}
\opt{D}{(III), (ii), (Q)}
\end{qoptions}
\qn{9} In which case particle will move with velocity changing periodically.
\begin{qoptions}
\opt{A}{(I), (iii), (R)}
\opt{B}{(IV), (i), (S)}
\opt{C}{(II), (iii), (S)}
\opt{D}{(III), (ii), (Q)}
\end{qoptions}
\qn{10} In which case particle will move on a straight line with constant speed.
\begin{qoptions}
\opt{A}{(I), (ii), (Q)}
\opt{B}{(I), (iii), (R)}
\opt{C}{(III), (ii), (S)}
\opt{D}{(IV), (i), (S)}
\end{qoptions}
\paraq{11}{13}
Column-1 represent distance of real object from optical instrument-1 given in Column-2. 2$^{\text{nd}}$ optical instrument place in contact with 1$^{\text{st}}$ is given column-3.\par
\begin{center}\footnotesize\setlength{\tabcolsep}{1pt}
\begin{tabular}{|L{0.7cm}L{0.8cm}|L{0.6cm}L{1.9cm}|L{0.5cm}L{1.9cm}|}
\hline
\multicolumn{2}{|c|}{Column-1}&\multicolumn{2}{c|}{Column-2}&\multicolumn{2}{c|}{Column-3}\\\hline
(I)&15 cm&(i)&\qfigi{p1_phy_s2_para11_13_1}&(P)&\qfigi{p1_phy_s2_para11_13_5}\\\hline
(II)&10 cm&(ii)&\qfigi{p1_phy_s2_para11_13_2}&(Q)&\qfigi{p1_phy_s2_para11_13_6}\\\hline
(III)&20 cm&(iii)&\qfigi{p1_phy_s2_para11_13_3}&(R)&\qfigi{p1_phy_s2_para11_13_7}\\\hline
(IV)&30 cm&(iv)&\qfigi{p1_phy_s2_para11_13_4}&(S)&\qfigi{p1_phy_s2_para11_13_8}\\\hline
\end{tabular}\end{center}
\qn{11} In which of following case final image will form on object itself.
\begin{qoptions}
\opt{A}{(II), (i), (Q)}
\opt{B}{(I), (i), (P)}
\opt{C}{(I), (i), (R)}
\opt{D}{(III), (iii), (R)}
\end{qoptions}
\qn{12} In which case distance between object and it's final virtual image will be 24 cm.
\begin{qoptions}
\opt{A}{(II), (iv), (P)}
\opt{B}{(III), (ii), (Q)}
\opt{C}{(IV), (iii), (S)}
\opt{D}{(I), (i), (R)}
\end{qoptions}
\qn{13} In which case real image of same size as object forms.
\begin{qoptions}
\opt{A}{(I), (ii), (R)}
\opt{B}{(II), (iii), (P)}
\opt{C}{(III), (iv), (R)}
\opt{D}{(IV), (ii), (S)}
\end{qoptions}
\sect{SECTION-II}
\qn{1} The top and bottom faces of a rigid insulating box of dimensions 1m $\times$ 1m $\times$ 1mm and mass 1 kg are charge with uniform surface charge densities $+10^{-2}$ C/m$^2$ and $-10^{-2}$ C/m$^2$ respectively and the box is placed in a uniform horizontal magnetic field $\vec B=(5\text{T})\hat j$. Assuming free space condition. Find velocity of box (in $\mu$m/s) when magnetic field is switched off.
\qfig{p1_phy_s3_q01_1}
\qn{2} Wind is blowing due east with speed 10 m/s. A car is moving west with speed 20 m/s towards a truck moving towards car with speed 10 m/s. If car blow horn of frequency 20 Hz, then find number of beats heard per second by a person in car. Assume waves reflected from truck reaches car. (speed of sound in air = 320 m/s)
\qn{3} A person can throw a ball to a maximum height of 20 m on the earth. If same person throw same ball on an imaginary planet whose density is twice of the earth then maximum possible radius of planet so that ball can escape from planet is $R$ km. Find $R$. (Radius of the Earth = $6.4\times10^6$ m)
\qn{4} A charged particle of charge 1$\mu$C and mass $10^{-6}$ kg is moving with a velocity of $9\times10^7\left(\hat i+2\hat k\right)$ m/s in space. There is a plane electromagnetic wave given by $\vec E=10^2\text{N/C}\,\hat i\sin(2\pi z-\omega t)$. The charge crosses origin at $t=\dfrac56$ns. At that moment, the magnitude of its acceleration is \underline{\hspace{1cm}} m/s$^2$ (Neglect relativistic effects)
\qn{5} A galvanometer is not connected to any circuit. When the needle is slightly deflected, it undergoes undamped oscillations with angular frequency of 5 rad/s. When the galvanometer terminals are connected by a conducting wire, it undergoes damped oscillations with angular frequency 4 rad/s. The radial magnetic field is 0.3T, area of coils is 100 cm$^2$ and there are 200 coils. If moment of inertia of assembly is $2.5\times10^{-3}$ kgm$^2$, estimate the resistance (in $\Omega$) of the circuit.
\qfig{p1_phy_s3_q05_1}\qfig{p1_phy_s3_q05_2}
\partsec{2}{CHEMISTRY}
\sect{SECTION-I(i)}
\qn{1} Which of the following is/are correct for the final product of the given sequence of reaction.
\qfig{p1_chem_s1_q01_1}
\begin{qoptions}
\opt{A}{Compound P$_3$ on oxidation with (CrO$_3$ + H$_2$SO$_4$) gives a compound which can give positive 2,4-DNP test.}
\opt{B}{Compound P$_3$ on reaction with I$_2$ + NaOH gives yellow precipitate.}
\opt{C}{Compound P$_3$ can react with PCC.}
\opt{D}{Compound P$_3$ on reaction with MnO$_2$ gives carboxylic acid.}
\end{qoptions}
\qn{2} Which of the following is/are correct statement regarding major product?
\qfig{p1_chem_s1_q02_1}
\begin{qoptions}
\opt{A}{Intramolecular nucleophilic substitution reaction takes place during the reaction sequence}
\opt{B}{5 membered ring is formed as major product}
\opt{C}{Grignard reagents (CH$_3$MgBr) prefer to give nucleophilic addition on polar $\pi$-bond}
\opt{D}{Cyclic ether is formed as final major product}
\end{qoptions}
\qn{3} For the following sequence, choose the correct option(s) ?
\qfig{p1_chem_s1_q03_1}
\begin{qoptions}
\opt{A}{(Q) will give positive Br$_2$, H$_2$O test}
\opt{B}{(Q) on reaction with neutral FeCl$_3$ gives violet colour}
\opt{C}{Major product (P) \& (S) are identical}
\opt{D}{Reaction sequence involves schmidt \& Hoffmann bromamide degradation reaction}
\end{qoptions}
\qn{4} Select correct statements -
\begin{qoptions}
\opt{A}{All adiabatic processes are isentropic}
\opt{B}{At equilibrium, at contant $T$ \& $P$, `$G$' becomes minimum}
\opt{C}{At equilibrium, entropy of universe becomes maximum}
\opt{D}{$G$ for a substance may change by changing $T$ \& $P$.}
\end{qoptions}
\qn{5} Four solutions of various electrolytes are given alongwith their percent ionization\par
(I) 0.1 M NaCl(aq.) (90 \% ionized)\par
(II) 0.05 M CaCl$_2$(aq.) (80 \% ionized)\par
(III) 0.04 M K$_3$[Fe(CN)$_6$](aq.) (60 \% ionized)\par
(IV) 0.03 M FeCl$_3$ (aq.) (70 \% ionized)\par
Identify correct statement(s) assuming all solutions to be dilute -
\begin{qoptions}
\opt{A}{Solution (IV) has maximum vapour pressure}
\opt{B}{Solution (I) has minimum vapour pressure}
\opt{C}{Solution (III) has lowest freezing point}
\opt{D}{Solution (I) has highest boiling point}
\end{qoptions}
\qn{6} In the following reactions. The correct statements is/are :\par
I. FeSO$_{4\,(\text{aq.})}$ + NO + H$_2$O $\to$ X\par
II. Sodium nitroprusside$_{(\text{aq.})}$ + Na$_2$S $\to$ Y
\begin{qoptions}
\opt{A}{Products X and Y both are paramagnetic}
\opt{B}{In the reaction I change of oxidation state of central atom occurs while in reaction II there is no change in oxidation state of metal}
\opt{C}{Product (X) can also be obtained when NO$_2^-$ or NO$_3^-$ react with dil. H$_2$SO$_4$ in presence of FeSO$_4$}
\opt{D}{Hybridisation of central atom of X is sp$^3$d$^2$ and Y is d$^2$sp$^3$}
\end{qoptions}
\qn{7} 1L of 0.2 M HCl (aq.) is mixed with 1 L of 0.2 M weak monoprotic acid HA ($K_a=5\times10^{-6}$) at $25^\circ$C.\par
Select the correct options :
\begin{qoptions}
\opt{A}{pH of the final solution will be 1.}
\opt{B}{\% dissociation of HA in final solution is approximately $5\times10^{-3}$ \%}
\opt{C}{Conc. of A$^-$ in final solution will be $1\times10^{-3}$ M}
\opt{D}{Conc. of A$^-$ in final solution will be $5\times10^{-6}$ M}
\end{qoptions}
\sect{SECTION-I(ii)}
\paraq{8}{10}
\begin{center}\footnotesize\setlength{\tabcolsep}{1pt}
\begin{tabular}{|L{0.75cm}L{2.0cm}|L{0.5cm}L{2.2cm}|L{0.5cm}L{1.8cm}|}
\hline
\multicolumn{2}{|c|}{Column - I}&\multicolumn{2}{c|}{Column - II}&\multicolumn{2}{c|}{Column - III}\\\hline
(I)&\qfigi{p1_chem_s2_pre_1}&(i)&Cold dilute alkaline KMnO$_4$&(P)&Meso compound\\\hline
(II)&\qfigi{p1_chem_s2_pre_2}&(ii)&Oxymercuration demercuration&(Q)&Threo compound is formed\\\hline
(III)&\qfigi{p1_chem_s2_pre_3}&(iii)&Hydroboration Oxidation&(R)&Racemic Mixture\\\hline
(IV)&\qfigi{p1_chem_s2_pre_4}&(iv)&Br$_2$/CCl$_4$&(S)&Diastereo Isomers\\\hline
\end{tabular}\end{center}
\qn{8} Which of the following combination is \textbf{True} ?
\begin{qoptions}
\opt{A}{III - iv - P}
\opt{B}{II - ii - Q}
\opt{C}{I - i - R}
\opt{D}{IV - i - R}
\end{qoptions}
\qn{9} Which of the following combination is \textbf{Incorrect} ?
\begin{qoptions}
\opt{A}{III - i - Q}
\opt{B}{III - i - R}
\opt{C}{IV - i - P}
\opt{D}{II - ii - R}
\end{qoptions}
\qn{10} Which of the following is a case of syn addition?
\begin{qoptions}
\opt{A}{III - i - R}
\opt{B}{III - ii - P}
\opt{C}{IV - iv - P}
\opt{D}{IV - iv - R}
\end{qoptions}
\columnbreak
\paraq{11}{13}
\begin{center}\footnotesize\setlength{\tabcolsep}{2pt}
\begin{tabular}{|L{2.2cm}|L{2.8cm}|L{2.3cm}|}
\hline
\multicolumn{1}{|L{2.2cm}|}{\centering Column-I\\ Formula of oxide of nitrogen}&\multicolumn{1}{L{2.8cm}|}{\centering Column-II\\ Common method of preparation}&\multicolumn{1}{L{2.3cm}|}{\centering Column-III\\ Chemical nature \& colour of oxide of nitrogen}\\\hline
(1) N$_2$O$_3$&(i) HNO$_3$ + P$_4$O$_{10}\to$&(P) neutral , colourless\\\hline
(2) NO$_2$&(ii) NO + N$_2$O$_4$ $\xrightarrow{-30^\circ\text{C}}$&(Q) acidic , blue\\\hline
(3) N$_2$O$_5$&(iii) Pb(NO$_3$)$_2$ $\xrightarrow{673\text{K}}$&(R) acidic , brown\\\hline
(4) N$_2$O&(iv) NH$_4$NO$_3$ $\xrightarrow{\Delta}$&(S) acidic , colourless\\\hline
\end{tabular}\end{center}
\qn{11} Which of the following is correctly matched?
\begin{qoptions}
\opt{A}{(1), (iv), (S)}
\opt{B}{(3), (iv), (P)}
\opt{C}{(4), (i), (R)}
\opt{D}{(2), (iii), (R)}
\end{qoptions}
\qn{12} Which of the following match is the CORRECT for N$_2$O$_5$?
\begin{qoptions}
\opt{A}{(i), (R)}
\opt{B}{(iv), (S)}
\opt{C}{(i), (S)}
\opt{D}{None of these}
\end{qoptions}
\qn{13} Which of the following is NOT correctly matched :-
\begin{qoptions}
\opt{A}{(1), (ii), (Q)}
\opt{B}{(4), (iii), (P)}
\opt{C}{(2), (iii), (R)}
\opt{D}{(4), (iv), (P)}
\end{qoptions}
\sect{SECTION-II}
\qn{1} For the reaction A(g) $\rightleftharpoons$ B(g) ; $K_p=10$ at 350K. Calculate $\Delta G$ (in kcal/mole) at 350 K, if partial pressures of A \& B gases are 10 bar \& 1 bar respectively. $\left[\text{Given : }R=2\dfrac{\text{cal}}{\text{mol}-\text{K}}\right]$
\qn{2} Using the data provided, calculate the multiple bond energy (kJ mol$^{-1}$) of a C$\equiv$C in C$_2$H$_2$. That energy is (take the bond energy of C--H bond as 350 kJ bond mol$^{-1}$).\par
2C(s) + H$_2$(g) $\to$ C$_2$H$_2$(g) \ ; $\Delta H=220$ kJ/mol\par
2C(s) $\to$ 2C(g) \ ; $\Delta H=1400$ kJ/mol\par
H$_2$(g) $\to$ 2H(g) \ ; $\Delta H=300$ kJ/mol\par
If the value of bond enthalpy is $x$ kJ/mol ; calculate $\left(\dfrac x{1000}\right)$
\qn{3}
\qfig{p1_chem_s3_q03_1}
Total number of anions which will give coloured gas/vapour will be :-\par
NO$_2^\ominus$, S$_2$O$_3^{2-}$, Cl$^-$, Br$^-$, I$^-$, NO$_3^-$, F$^-$
\qn{4} Aqueous CuSO$_4$ decolorizes on addition of KCN (excess) due to formation of complex (A).\par
In complex ``A''. If\par
(i) Number of d orbitals participating in hybridisation is / are `a'\par
(ii) Coordination number of Cu is ``c''.\par
Then find 8a + 5c.
\qn{5} Mg crystallizes in HCP. Calculate total number of tetrahedral and octahedral voids in 12 milligrams of Mg ($N_A=6\times10^{23}$) If your answer is $x\times10^{20}$ then fill value of $x$ in OMR sheet.
\partsec{3}{MATHEMATICS}
\sect{SECTION-I(i)}
\qn{1} Let $f$, $g$ and $h$ be function differentiable on some open interval around 0 and satisfy the equation $f'=2f^2gh+\dfrac1{gh}$, $f(0)=1$, $g'=fg^2h+\dfrac4{fh}$, $g(0)=1$ and $h'=3fgh^2+\dfrac1{fg}$, $h(0)=1$. The function $f$ is given by
\begin{qoptions}
\opt{A}{$f(x)=2^{-1/12}\left(\dfrac{\sin\left(6x+\frac\pi4\right)}{\cos^2\left(6x+\frac\pi4\right)}\right)^{1/6}$}
\opt{B}{$f(x)=2^{-1/6}\left(\dfrac{\sin\left(6x+\frac\pi4\right)}{\cos^2\left(6x+\frac\pi4\right)}\right)^{1/12}$}
\opt{C}{$f(x)=(\sec12x)^{1/12}(\sec12x+\tan12x)^{1/4}$}
\opt{D}{$f(x)=(\sec12x)^{1/4}(\sec12x+\tan12x)^{1/12}$}
\end{qoptions}
\qn{2} A differentiable function $f:R\to R$ satisfies following relations\par
(i) $f(x)-f(y)=(x-y)f'\left(\dfrac{x+y}2\right)$, $\forall x,y\in R$\par
(ii) $f(0)=f(1)=0$, $f(2)>0$\par
Which of the following is/are true
\begin{qoptions}
\opt{A}{$f(x)$ is an onto function}
\opt{B}{$f(x)$ has an extrema at $x=\dfrac12$}
\opt{C}{$f(x)$ is a many one function}
\opt{D}{$\displaystyle\lim_{x\to-\infty}f(x)=-\infty$}
\end{qoptions}
\qn{3} Let $A$ be the set of $2\times2$ matrices given by $A=[a_{ij}]$; where $a_{ij}\in\{0,1,2,3,4,5\}$ such that $a_{11}+a_{12}+a_{21}+a_{22}=5$, then
\begin{qoptions}
\opt{A}{the number of matrices in $A$ which are invertible are 32}
\opt{B}{the number of matrices in $A$ which are symmetric or skew symmetric are 8}
\opt{C}{the number of matrices in $A$ which are symmetric or skew symmetric such that Det $A$ is odd are 6}
\opt{D}{The total number of matrices in $A$ are 56}
\end{qoptions}
\qn{4} $\vec a$ and $\vec b$ be two vectors such that $|\vec a|=1$, $|\vec b|=4$ and $\vec a\cdot\vec b=2$. If $\vec c=\left(2\vec a\times\vec b\right)-3\vec b$, then which of the following is/are correct ?
\begin{qoptions}
\opt{A}{$\vec b\cdot\vec c=48$}
\opt{B}{$\vec b\cdot\vec c=-48$}
\opt{C}{Angle between $\vec b$ and $\vec c$ is $\dfrac{5\pi}6$}
\opt{D}{Angle between $\vec b$ and $\vec c$ is $\dfrac\pi6$}
\end{qoptions}
\qn{5} Consider a pyramid built on a square base of side 2 such that the lengths of the line segments joining its apex to the 4 corners of the base are equal and the height of the structure is equal to the length of a side of its base. Then which of the following statements is/are true
\begin{qoptions}
\opt{A}{volume of the pyramid is $\dfrac83$}
\opt{B}{the angle between the base and any other face is $\sin^{-1}\dfrac2{\sqrt5}$}
\opt{C}{the angle between two faces with common edge which passes through the apex is $\cos^{-1}\dfrac15$}
\opt{D}{volume of the largest tetrahedron that can be inscribed in the pyramid is 1}
\end{qoptions}
\qn{6} If volume of a parallelepiped determined by the vectors $\vec a,\vec b$ and $\vec c$ is 2, then volume of parallelepiped determined by the vectors
\begin{qoptions}
\opt{A}{$3(\vec a+\vec b)$, $(\vec b+\vec c)$ and $(\vec c+\vec a)$ is 24}
\opt{B}{$2(\vec a\times\vec b)$, $3(\vec b\times\vec c)$ and $(\vec c\times\vec a)$ is 24}
\opt{C}{$(\vec a-2\vec b)$, $(\vec b-2\vec c)$ and $(\vec c-2\vec a)$ is 14}
\opt{D}{$(3\vec a-\vec b)$, $(3\vec b-\vec c)$ and $(3\vec c-\vec a)$ is 52}
\end{qoptions}
\qn{7} Let $(1+\sqrt5)^n=\alpha_n+\beta_n\sqrt5$ where $\alpha_n$ and $\beta_n$ are integers, then choose the correct option(s)
\begin{qoptions}
\opt{A}{$\displaystyle\lim_{\theta\to0}\frac{\tan^{-1}\left(\frac{\alpha_{n+1}-\beta_{n+1}}3\right)\theta}{\sin^{-1}(\beta_n\theta)}=\dfrac43$}
\opt{B}{$\displaystyle\lim_{\theta\to0}\frac{\tan^{-1}\left(\frac{\alpha_{n-1}-\beta_{n+1}}3\right)\theta}{\sin^{-1}(\beta_n\theta)}=\dfrac34$}
\opt{C}{$\left(\dfrac{\alpha_n}{2^n}\right)^2-5\left(\dfrac{\beta_n}{2^n}\right)^2=(-1)^n$}
\opt{D}{$\left(\dfrac{\alpha_n}{2^n}\right)^2-5\left(\dfrac{\beta_n}{2^n}\right)^2=(-1)^{\frac n2}$}
\end{qoptions}
\sect{SECTION-I(ii)}
\paraq{8}{10}
\noindent\textbf{Column-I} : $y=f(x)$ with respective domain.\par
\noindent\textbf{Column-II} : Interval in which $f(x)$ is increasing\par
\noindent\textbf{Column-III} : Interval in which $f(x)$ is decreasing\par
\begin{center}\footnotesize\setlength{\tabcolsep}{2pt}
\begin{tabular}{|L{3.6cm}|L{1.9cm}|L{1.9cm}|}
\hline
\multicolumn{1}{|c|}{Column-I}&\multicolumn{1}{c|}{Column-II}&\multicolumn{1}{c|}{Column-III}\\\hline
(I) $y=\sin x+\cos x$ in $\left(0,\frac\pi2\right)$&(i) $\left(0,\frac\pi4\right)$&(P) $\left(0,\frac\pi4\right)$\\\hline
(II) $y=\sin2x$ in $\left(0,\frac\pi2\right)$&(ii) $\left(\frac\pi4,\frac\pi2\right)$&(Q) $\left(\frac\pi4,\frac\pi2\right)$\\\hline
(III) $y=\sin^5x+\cos^5x-1$ in $\left(0,\frac\pi2\right)$&(iii) $\left(0,\frac\pi3\right)$&(R) $\left(0,\frac1e\right)$\\\hline
(IV) $y=x^x$ in $\left(0,\frac\pi2\right)$&(iv) $\left(\frac1e,\frac\pi2\right)$&(S) $\left(\frac1e,\frac\pi2\right)$\\\hline
\end{tabular}\end{center}
\qn{8} Which of the following options is the only \textbf{CORRECT} combination ?
\begin{qoptions}
\opt{A}{(I) (i) (S)}
\opt{B}{(II) (iii) (S)}
\opt{C}{(III) (ii) (P)}
\opt{D}{(I) (ii) (P)}
\end{qoptions}
\qn{9} Which of the following options is the only \textbf{CORRECT} combination ?
\begin{qoptions}
\opt{A}{(I) (iii) (P)}
\opt{B}{(II) (ii) (P)}
\opt{C}{(III) (iii) (P)}
\opt{D}{(IV) (iv) (R)}
\end{qoptions}
\qn{10} Which of the following options is the only \textbf{CORRECT} combination ?
\begin{qoptions}
\opt{A}{(I) (iv) (R)}
\opt{B}{(II) (i) (Q)}
\opt{C}{(III) (iii) (Q)}
\opt{D}{(II) (ii) (P)}
\end{qoptions}
\paraq{11}{13}
Answer Q.11, Q.12 and Q.13 by appropriately matching the information given in the three columns of the following table\par
\begin{center}\footnotesize\setlength{\tabcolsep}{1pt}
\begin{tabular}{|L{3.5cm}|L{2.8cm}|L{1.2cm}|}
\hline
\multicolumn{1}{|c|}{Column 1}&\multicolumn{1}{c|}{Column 2}&\multicolumn{1}{c|}{Column 3}\\\hline
(I) $\displaystyle\int_0^{\pi/2}\sin^8x\,dx=\dfrac{A\pi}{2^B}$ (where $A$ and $2^B$ are co-prime numbers), then $\dfrac{A}{B-1}$ is&(i) If slope of tangent to the curve $y=x+\sin x\cos x$ at $x=\dfrac\pi2$ is $m$, then $(m+7)$ is&(P) odd\\\hline
(II) $\dfrac{\displaystyle\int_0^8[x]\,dx}{\displaystyle\int_0^8\{x\}\,dx}$ is (where $[.]$ is greatest integer function and $\{.\}$ fractional part function respectively)&(ii) If $y$-intercept of perpendicular bisector of $P(2,5)$ and $Q(k,2)$ is $\dfrac32$, then $k(k>0)$ can be&(Q) even\\\hline
(III) If $\ell=\displaystyle\lim_{t\to0}\frac1{t^4}\int_0^t\frac{x^2\ell n(1+x)}{x^7+3}\,dx$ then $\dfrac3{4\ell}$ is&(iii) Area of polygon formed by $x^2+x-2=0$ and $y^2-5y+4=0$ is&(R) prime\\\hline
(IV) $\displaystyle\lim_{n\to\infty}\Bigl(\dfrac1{\sqrt n}+\dfrac1{\sqrt{2n}}+\dfrac1{\sqrt{3n}}+\dots+\dfrac1{2n}\Bigr)$ is&(iv) Rate of change of area of an equilateral triangle w.r.t. its side when side is $\dfrac{10}{\sqrt3}$ is&(S) composite\\\hline
\end{tabular}\end{center}
\qn{11} Which of the following options is the only \textbf{CORRECT} combination ?
\begin{qoptions}
\opt{A}{(I) (iii) (Q,R)}
\opt{B}{(I) (i) (P,R)}
\opt{C}{(I) (iv) (P,R)}
\opt{D}{(I) (ii) (Q,R)}
\end{qoptions}
\qn{12} Which of the following options is the only \textbf{CORRECT} combination ?
\begin{qoptions}
\opt{A}{(II) (i) (P,R)}
\opt{B}{(II) (ii) (P,R)}
\opt{C}{(II) (iii) (Q,S)}
\opt{D}{(II) (iv) (Q,S)}
\end{qoptions}
\qn{13} Which of the following options is the only \textbf{CORRECT} combination ?
\begin{qoptions}
\opt{A}{(III) (ii) (P,S)}
\opt{B}{(III) (i) (Q,S}
\opt{C}{(IV) (ii) (Q,S)}
\opt{D}{(IV) (iii) (P,S)}
\end{qoptions}
\sect{SECTION-II}
\qn{1} If in a $4\times4$ matrix the sum of elements of each row, column and both the main diagonals is 1, then sum of the four corner elements is equal to
\qn{2} Let $A=\begin{bmatrix}5&17\\13&29\end{bmatrix}$ and $B$ is $2\times2$ matrix whose elements are elements of set $P\{3,5,7,11,13,17,19,23,29\}$.\par
If all the elements of $AB$ are multiple of 4, then the total number of possible matrices $B$ with distinct elements.
\qn{3} Consider a $3\times3$ matrix $A=\{a_{ij}\}$ where $a_{ij}=\dfrac1{i+j-1}\ \forall\ i,j\in\{1,2,3\}$. If matrix $B=\{b_{ij}\}$ of order $3\times3$ exists such that $A^{-1}=3B$, then the value of $\dfrac{|B|}{a_{23}\cdot b_{22}}$ is (where $|X|$ denotes determinant value of matrix $X$)
\qn{4} In a $\Delta ABC$, $AB=4$, $AC=5$, $O$ is circumcenter of $\Delta ABC$, then the value of $2\overrightarrow{AO}\cdot\overrightarrow{BC}$ ?
\qn{5} If $a_n=\left(\dfrac{2n+1}{n+1}\right){}^8C_n$; $n=\{0,1,2,\dots,7,8\}$\par
$b_n=\left(\dfrac{3n+1}{n+1}\right){}^7C_n$; $n=\{0,1,2,\dots\dots,7\}$\par
If $a_p=\max\limits_{0\le n\le8}(a_n)$ and $b_q=\min\limits_{0\le n\le7}(b_n)$, where $n\in W$ then the value of $(p+q)$ is
\raggedcolumns
\papersec{2}
\partsec{1}{PHYSICS}
\sect{SECTION-I}
\qn{1} A 750-nm light beam hits the flat surface of a certain liquid, and the beam is split into a reflected ray and a refracted ray. If the reflected ray is completely polarized and angle of refraction is $37^\circ$, what is the wavelength of the refracted ray ?
\begin{qoptions}
\opt{A}{1000 nm}
\opt{B}{562.5 nm}
\opt{C}{500 nm}
\opt{D}{1125 nm}
\end{qoptions}
\qn{2} The path difference between two interfering waves at a point on the screen is $\lambda/6$. The ratio of intensity at this point and that at the central bright fringe will be : (Assume that intensity due to each slit in same)
\begin{qoptions}
\opt{A}{0.853}
\opt{B}{8.53}
\opt{C}{0.75}
\opt{D}{7.5}
\end{qoptions}
\qn{3} If we study the vibration of a pipe open at both ends, then the following statements is not true :
\begin{qoptions}
\opt{A}{Only odd harmonics of the fundamental frequency will be generated.}
\opt{B}{All harmonics of the fundamental frequency will be generated.}
\opt{C}{Open end will be pressure node}
\opt{D}{Open end will be displacement antinode}
\end{qoptions}
\qn{4} A standing wave having 3 nodes and 2 antinodes is formed between two atoms having a distance of 1.21 \AA\ between them. The wavelength of the standing wave is -
\begin{qoptions}
\opt{A}{1.21 \AA}
\opt{B}{2.42 \AA}
\opt{C}{3.63 \AA}
\opt{D}{6.05 \AA}
\end{qoptions}
\qn{5} A closed organ pipe and an open organ pipe are tuned to the same fundamental frequency. What is the ratio of their lengths ?
\begin{qoptions}
\opt{A}{1 : 2}
\opt{B}{2 : 1}
\opt{C}{2 : 3}
\opt{D}{4 : 3}
\end{qoptions}
\qn{6} A traveller is running with velocity $v_0$ towards a stationary train. If the train sounds a horn of frequency $n$ then the apparent frequency heard by the traveller will be ($v$ is speed of sound) :
\begin{qoptions}
\opt{A}{$n'=n\left(\dfrac{v+v_0}{v}\right)$}
\opt{B}{$n'=n\left(\dfrac{v-v_0}{v+v_s}\right)$}
\opt{C}{$n'=n\left(\dfrac{v-v_0}{v}\right)$}
\opt{D}{None of these}
\end{qoptions}
\qn{7} The equation of a progressive wave is $y=0.02\sin2\pi\left[\dfrac t{0.01}-\dfrac x{0.3}\right]$, where $x$ and $y$ are in meters and $t$ is in second. The velocity (in m/s) of propagation of the wave is :
\begin{qoptions}
\opt{A}{400}
\opt{B}{40}
\opt{C}{300}
\opt{D}{30}
\end{qoptions}
\qn{8} In a sinusoidal wave, the time required for a particular point to move from maximum displacement to zero displacement is 0.17 sec. The frequency of the wave is :-
\begin{qoptions}
\opt{A}{1.47 Hz}
\opt{B}{0.36 Hz}
\opt{C}{0.73 Hz}
\opt{D}{2.94 Hz}
\end{qoptions}
\qn{9} Laser light with a wavelength of 500 nm is directed through one slit or two slits and allowed to fall on a screen 2.60 m beyond. Figure shows the pattern on the screen, with a centimeter ruler below it. Did the light pass through one slit or two slits ? If one slit, find its width. If two slits, find the distance between their centers.
\qfig{p2_phy_s1_q09_1}
\begin{qoptions}
\opt{A}{One slit of width approx 0.1 mm.}
\opt{B}{One slit of width approx 0.05 mm.}
\opt{C}{Two slits at an approx. distance of 0.08 mm}
\opt{D}{Two slits at an approx. distance of 0.04 mm}
\end{qoptions}
\qn{10} The antenna of a 6 kW radio station radiates spherical wavefronts. At a distance of 100 m from the antenna, what is the peak magnetic field strength ?
\begin{qoptions}
\opt{A}{20 nT}
\opt{B}{$10\sqrt2$nT}
\opt{C}{10nT}
\opt{D}{$\dfrac5{\sqrt2}$nT}
\end{qoptions}
\qn{11} A particle executes SHM and is located at $x=a$, $x=b$ and $x=c$, at $t=t_0$, $t=2t_0$ and $t=3t_0$ respectively. The angular frequency of oscillation is
\begin{qoptions}
\opt{A}{$\dfrac1{t_0}\sin^{-1}\left(\dfrac{a+c}{2b}\right)$}
\opt{B}{$\dfrac1{t_0}\sin^{-1}\left(\dfrac{a+c}{b}\right)$}
\opt{C}{$\dfrac1{t_0}\cos^{-1}\left(\dfrac{a+c}{b}\right)$}
\opt{D}{$\dfrac1{t_0}\cos^{-1}\left(\dfrac{a+c}{2b}\right)$}
\end{qoptions}
\qn{12} A uniform hemispherical shell of mass $M$ and radius $R$ is made to perform small angular oscillation about point `$A$' such that line $AO$ always remains in vertical plane (where $O$ is curvature centre of spherical shell). Frequency of small angular oscillations will be :
\qfig{p2_phy_s1_q12_1}
\begin{qoptions}
\opt{A}{$\dfrac1{2\pi}\sqrt{\dfrac{6g}R}$}
\opt{B}{$\dfrac1{2\pi}\sqrt{\dfrac{3g}R}$}
\opt{C}{$\dfrac1{4\pi}\sqrt{\dfrac{3g}R}$}
\opt{D}{$\dfrac1{4\pi}\sqrt{\dfrac gR}$}
\end{qoptions}
\qn{13} A non-uniform rope of length $\ell$ hangs from a ceilling. Mass per unit length of rope ($\mu$) changes as $\mu=\mu_0e^y$, where $y$ is the distance along the string from its lowest point. Then graph between square of velocity of wave $\left(v_y^2\right)$ and $y$ will be best represented as :
\qfig{p2_phy_s1_q13_1}
\begin{qoptions}
\opt{A}{\qfigi{p2_phy_s1_q13_A}}
\opt{B}{\qfigi{p2_phy_s1_q13_B}}
\opt{C}{\qfigi{p2_phy_s1_q13_C}}
\opt{D}{None of these}
\end{qoptions}
\qn{14} $S_1$ and $S_2$ are two coherent sources of radiations separated by distance $100.25\ \lambda$, where $\lambda$ is the wave length of radiation. $S_1$ leads $S_2$ in phase by $\dfrac\pi2$. $A$ and $B$ are two points on the line joining $S_1$ and $S_2$ as shown in figure. The ratio of amplitudes of component waves from source $S_1$ and $S_2$ at $A$ and $B$ are in ratio $1:2$. The ratio of intensity at $A$ to that of $B\left(\dfrac{I_A}{I_B}\right)$ is :
\qfig{p2_phy_s1_q14_1}
\begin{qoptions}
\opt{A}{$\infty$}
\opt{B}{$\dfrac19$}
\opt{C}{0}
\opt{D}{8.1}
\end{qoptions}
\qn{15} In the given circuit key is closed at $t=0$. At what time the P.D. across inductor is one fourth of emf of the cell.
\qfig{p2_phy_s1_q15_1}
\begin{qoptions}
\opt{A}{$\dfrac LR\ell n2$}
\opt{B}{$\dfrac LR\ell n4$}
\opt{C}{0}
\opt{D}{$\dfrac{2L}R$}
\end{qoptions}
\qn{16} In the ac circuit shown, $X_L=7\Omega$, $R=4\Omega$ and $X_C=4\Omega$. The reading of the ideal voltmeter $V_2$ is $8\sqrt2$ V. The reading of the ideal ammeter is.
\qfig{p2_phy_s1_q16_1}
\begin{qoptions}
\opt{A}{3A}
\opt{B}{2A}
\opt{C}{4A}
\opt{D}{5A}
\end{qoptions}
\qn{17} The arrangements of wires in which there exists atleast one point at finite distance from wires, at which magnetic field is zero.
\begin{qoptions}
\opt{A}{\qfigi{p2_phy_s1_q17_A}}
\opt{B}{\qfigi{p2_phy_s1_q17_B}}
\opt{C}{\qfigi{p2_phy_s1_q17_C}}
\opt{D}{\qfigi{p2_phy_s1_q17_D}}
\end{qoptions}
\qn{18} If the earth's field induction at a place is 0.36 gauss and the angle of dip is $60^\circ$ then the horizontal and vertical component of the field is :
\begin{qoptions}
\opt{A}{0.36 gauss, $0.36\sqrt3$ gauss}
\opt{B}{0.18 gauss, $0.18\sqrt3$gauss}
\opt{C}{0.09 gauss, $0.09\sqrt3$ gauss}
\opt{D}{None of these}
\end{qoptions}
\qn{19} In a series L, C circuit, which of the following represents variation of magnitude of reactance ($X$) with frequency ($f$)
\begin{qoptions}
\opt{A}{\qfigi{p2_phy_s1_q19_A}}
\opt{B}{\qfigi{p2_phy_s1_q19_B}}
\opt{C}{\qfigi{p2_phy_s1_q19_C}}
\opt{D}{\qfigi{p2_phy_s1_q19_D}}
\end{qoptions}
\qn{20} Three copper wires of length and cross-sectional area $(L,A)$, $(2L,A/2)$ and $(L/2,2A)$. Resistance is minimum in
\begin{qoptions}
\opt{A}{wire of cross-sectional area $A$}
\opt{B}{wire of area $A/2$}
\opt{C}{wire of cross-sectional area $2A$}
\opt{D}{same in all three cases.}
\end{qoptions}
\sect{SECTION-II}
\qn{1} Refractive indices of 2 different media with separating boundary at the diagonal of rectangular glass slab are shown. Then angle of deviation (in degree) of the ray in going from $n_1$ to $n_2$ as shown in the figure is :
\qfig{p2_phy_s2_q01_1}
\qn{2} A planet is made of two materials of density `$\rho_1$' and `$\rho_2$' as shown in figure. Acceleration due to gravity at surface of planet is same as that at depth `$R$'. The ratio $\dfrac{\rho_1}{\rho_2}$ is equal to $\dfrac x3$. Find the value of `$x$'.
\qfig{p2_phy_s2_q02_1}
\qn{3} The potential energy of a particle of mass $m$ free to move along $x$ axis has the form $U(x)=\dfrac a{x^2}-\dfrac bx$, where $a$ and $b$ are positive constant and $x$ is $x$-coordinate of the particle. The time period of SHM (along $x$-axis) of small amplitude around mean position is $\dfrac{P\pi a}{b^2}\sqrt{2ma}$. Find the value of $P$.
\qn{4} Radii of curvature of a concavo-convex lens (refractive index = 1.5) are 40 cm (concave side) and 20 cm (convex side) as shown. The convex side is silvered. The distance $x$ on the principal axis where object is placed so that its image is created on the object itself, is equal to (in cm) :
\qfig{p2_phy_s2_q04_1}
\qn{5} The frequency of a fork A is 3\% more than the frequency of a standard fork whereas the frequency of fork B is 3\% less. The forks A and B produce 6 beats per second. The frequency of standard fork (in Hz) will be :-
\qn{6} If the ratio of the intensity of two coherent sources is 4 then the visibility $\dfrac{I_{\max}-I_{\min}}{I_{\max}+I_{\min}}$ of the fringes is $n$. Find the value of $10n$.
\qn{7} Figure shows two coherent sources $S_1$ and $S_2$ vibrating in same phase, $AB$ is an wire lying symmetrically at a far distance from the sources $S_1$ and $S_2$. Let $\dfrac\lambda d=10^{-3}$. $\angle BOA=0.12^\circ$. How many bright spots will be seen on the wire, including points $A$ and $B$.
\qfig{p2_phy_s2_q07_1}
\qn{8} A graph of the $x$ component of the electric field as a function of $x$ in a region of space is shown. The $Y$ and $Z$ components of the electric field are zero in this region. If the electric potential is 10V at the origin, then potential (in V) at $x=2.0$ m is :
\qfig{p2_phy_s2_q08_1}
\qn{9} Consider a uniformly charged solid non conducting sphere of radius $R$ and charge $Q$. Consider three spherical regions centered at the centre of the uniformly charged solid non conducting sphere, as given below :\par
Region (1) $R_1\ge r\ge0$\par
Region (2) $R_2\ge r\ge R_1$\par
Region (3) $\infty\ge r\ge R_2$\par
(Both $R_1$ \& $R_2>R$)\par
If electrostatic potential energy stored in each of the three regions is equal then $\dfrac{R_2}{R_1}$ is :
\qn{10} A sound absorber attenuates the sound level by 20 dB. The intensity decreases by a factor of :
\partsec{2}{CHEMISTRY}
\sect{SECTION-I}
\qn{1} Select the incorrect statement about three-dimensional HCP crystal made of identical atoms : [Given : $r$ = radius of atom]
\begin{qoptions}
\opt{A}{The number of atoms in HCP unit cell is six}
\opt{B}{The volume of HCP unit cell is $24\sqrt2\,r^3$}
\opt{C}{The empty space in HCP unit cell is 26\%}
\opt{D}{The base area of HCP unit cell is $1.5\sqrt3\,r^2$}
\end{qoptions}
\qn{2} For a reaction at equilibrium
\par\noindent\hfill$\mathrm A(g)\rightleftharpoons \mathrm B(g)+\tfrac12\mathrm C(g)$\hfill\null\par
the relation between dissociation constant ($K$), degree of dissociation ($\alpha$) and equilibrium pressure ($p$) is given by :
\begin{qoptions}
\opt{A}{$K=\dfrac{\alpha^{1/2}p^{3/2}}{\left(1+\frac32\alpha\right)^{1/2}(1-\alpha)}$}
\opt{B}{$K=\dfrac{\alpha^{3/2}p^{1/2}}{(2+\alpha)^{1/2}(1-\alpha)}$}
\opt{C}{$K=\dfrac{(\alpha p)^{3/2}}{\left(1+\frac32\alpha\right)^{1/2}(1-\alpha)}$}
\opt{D}{$K=\dfrac{(\alpha p)^{3/2}}{(1+\alpha)(1-\alpha)^{1/2}}$}
\end{qoptions}
\qn{3} Enthalpy of neutralization of H$_3$PO$_3$ acid is $-106.68$ kJ/mol using NaOH. Enthalpy of neutralization of HCl by NaOH is $-55.84$ kJ/mol. Calculate $\Delta H_{\text{ionization}}$ of H$_3$PO$_3$ into its ions:
\begin{qoptions}
\opt{A}{50.84 kJ/mol}
\opt{B}{5 kJ/mol}
\opt{C}{2.5 kJ/mol}
\opt{D}{None of these}
\end{qoptions}
\qn{4} The enthalpy of combustion of propane (C$_3$H$_8$) gas in terms of given data is :
\par\smallskip\noindent Bond energy (kJ/mol)\par\smallskip
\begin{center}\footnotesize\setlength{\tabcolsep}{2pt}
\begin{tabular}{|c|c|c|c|c|}\hline
$\varepsilon_{\text{C--H}}$&$\varepsilon_{\text{O=O}}$&$\varepsilon_{\text{C=O}}$&$\varepsilon_{\text{O--H}}$&$\varepsilon_{\text{C--C}}$\\\hline
$+x_1$&$+x_2$&$+x_3$&$+x_4$&$+x_5$\\\hline
\end{tabular}\end{center}
Resonance energy of CO$_2$ is $-\,z$ kJ/mol and $\Delta H_{\text{vaporization}}$ [H$_2$O($\ell$)] is $y$ kJ/mol.
\begin{qoptions}
\opt{A}{$8x_1+2x_5+5x_2-6x_3-8x_4-4y-3z$}
\opt{B}{$6x_1+x_5+5x_2-3x_3-4x_4-4y-3z$}
\opt{C}{$8x_1+2x_5+5x_2-6x_3-8x_4-y-z$}
\opt{D}{$8x_1+x_5+5x_2-6x_3-8x_4-4y+3z$}
\end{qoptions}
\qn{5} What will be the order of reaction [A $\to$ P] and rate constant for this chemical change having [$\log t_{50\%}$] vs [$\log$ initial concentration of A (a)] curves as :
\qfig{p2_chem_s1_q05_1}
\begin{qoptions}
\opt{A}{0, 1/2}
\opt{B}{1, 1}
\opt{C}{2, 2}
\opt{D}{3, 1}
\end{qoptions}
\qn{6} \qfig{p2_chem_s1_q06_1}
Major product of the reaction is:
\begin{qoptions}
\opt{A}{\qfigi{p2_chem_s1_q06_A}}
\opt{B}{\qfigi{p2_chem_s1_q06_B}}
\opt{C}{\qfigi{p2_chem_s1_q06_C}}
\opt{D}{\qfigi{p2_chem_s1_q06_D}}
\end{qoptions}
\qn{7} Which of the following Amine can be synthesised in good yield by Gabriel phthalimide synthesis?
\begin{qoptions}
\opt{A}{Ph--NH$_2$}
\opt{B}{CH$_3$--N(CH$_3$)--CH$_3$}
\opt{C}{Ph--NH--CH$_3$}
\opt{D}{\qfigi{p2_chem_s1_q07_D}}
\end{qoptions}
\qn{8} Which of the following is incorrectly matched with their product?
\begin{qoptions}
\opt{A}{$2\,\mathrm{CH_3{-}CO{-}CH_3}\xrightarrow{\ominus\mathrm{OH}/\Delta}$ mesityl oxide}
\opt{B}{\qfigi{p2_chem_s1_q08_2}\ $\xrightarrow{\mathrm{CHCl_3+NaOH}}$ Salicylaldehyde}
\opt{C}{\qfigi{p2_chem_s1_q08_C}}
\opt{D}{\qfigi{p2_chem_s1_q08_D}}
\end{qoptions}
\qn{9} \par\noindent\hfill cyclohexene $\xrightarrow{\mathrm{CH_3CO_3H}}$ Product\hfill\null\par
Product formed in the reaction is -
\begin{qoptions}
\opt{A}{\qfigi{p2_chem_s1_q09_A}}
\opt{B}{\qfigi{p2_chem_s1_q09_B}}
\opt{C}{\qfigi{p2_chem_s1_q09_C}}
\opt{D}{\qfigi{p2_chem_s1_q09_D}}
\end{qoptions}
\qn{10} Which of the following is supposed to produce alkene
\begin{qoptions}
\opt{A}{CH$_3$--Cl $\xrightarrow[\text{D.E.}]{\text{Na}}$}
\opt{B}{\qfigi{p2_chem_s1_q10_B}}
\opt{C}{CH$_3$CH$_2$--OH $\xrightarrow{\text{PCC}}$}
\opt{D}{CH$_3$CH$_2$--CN $\xrightarrow[\text{H}_2\text{O}]{\text{SnCl}_2+\text{HCl}}$}
\end{qoptions}
\qn{11} \par\noindent\hfill$\mathrm{CH_3{-}CH_2{-}CH{=}O+\underset{\text{(excess)}}{HCHO}}\xrightarrow{\ \overline{\mathrm{O}}\mathrm{H}\ }$\hfill\null\par
Which of the following product mixture is formed as a final major product
\begin{qoptions}
\opt{A}{\qfigi{p2_chem_s1_q11_A}}
\opt{B}{\qfigi{p2_chem_s1_q11_B}}
\opt{C}{\qfigi{p2_chem_s1_q11_C}}
\opt{D}{\qfigi{p2_chem_s1_q11_D}}
\end{qoptions}
\qn{12} In the following reaction, the compound J is :
\qfig{p2_chem_s1_q12_1}
J gives effervescence on treatment with NaHCO$_3$ and a positive Baeyer's test.

The compound J is :
\begin{qoptions}
\opt{A}{\qfigi{p2_chem_s1_q12_A}}
\opt{B}{\qfigi{p2_chem_s1_q12_B}}
\opt{C}{\qfigi{p2_chem_s1_q12_C}}
\opt{D}{\qfigi{p2_chem_s1_q12_D}}
\end{qoptions}
\qn{13} \qfig{p2_chem_s1_q13_1}
Final product Y is :
\begin{qoptions}
\opt{A}{O=CH--CH$_2$--CH$_2$--CH$_2$--CH$_2$--CH=O}
\opt{B}{\qfigi{p2_chem_s1_q13_B}}
\opt{C}{\qfigi{p2_chem_s1_q13_C}}
\opt{D}{\qfigi{p2_chem_s1_q13_D}}
\end{qoptions}
\qn{14} Choose \textbf{CORRECT} statement(s) regarding borax :
\par\noindent(i) Its aqueous solution is acidic in nature\par
\noindent(ii) It is prepared by using colemanite mineral\par
\noindent(iii) Its aqueous solution acts as buffer solution\par
\noindent(iv) It produces brick red transparent bead on heating with CaCO$_3$
\begin{qoptions}
\opt{A}{(i), (ii)(iii) only}
\opt{B}{(ii), (iii) only}
\opt{C}{(ii), (iii)(iv) only}
\opt{D}{All of these}
\end{qoptions}
\qn{15} Which of the following statement is correct for H$_2$O and D$_2$O?
\begin{qoptions}
\opt{A}{D$_2$O has higher number of valence shell electron than H$_2$O}
\opt{B}{D$_2$O has lower value of boiling point than H$_2$O}
\opt{C}{D$_2$O has higher number of neutrons than H$_2$O}
\opt{D}{H$_2$O has smaller number of protons than D$_2$O}
\end{qoptions}
\qn{16} Choose incorrect order regarding mentioned properties :
\begin{qoptions}
\opt{A}{SeO$_2$ $>$ SO$_2$ (Reducing power)}
\opt{B}{Na $>$ K (Softness of metal)}
\opt{C}{PbCl$_4$ $>$ PbI$_4$ (stability)}
\opt{D}{H$_2$S $>$ H$_2$O (Acidic character)}
\end{qoptions}
\qn{17} Which of the reaction does NOT give product as indicated ? (ignore balancing)
\begin{qoptions}
\opt{A}{PbS(s) + O$_3$(g) $\to$ PbSO$_4$(s) + O$_2$(g)}
\opt{B}{ZnS(s) + O$_2$(g) $\to$ ZnO(s) + SO$_2$(g)}
\opt{C}{Pb$_2$O$_3$ + HNO$_3$ $\to$ Pb(NO$_3$)$_2$ + PbO + H$_2$O}
\opt{D}{KClO$_3$(s) $\xrightarrow{\text{MnO}_2}$ KCl(s) + O$_2$(g)}
\end{qoptions}
\qn{18} Among statements (a)--(d), the correct ones are :
\par\noindent(a) Lime stone is decomposed to CaO and combined with SiO$_2$ during the extraction of iron from its oxides.\par
\noindent(b) In the extraction of silver, silver is leached out as an anionic complex.\par
\noindent(c) Nickel is purified by Mond's process.\par
\noindent(d) Zr and Ti are purified by Van Arkel method.
\begin{qoptions}
\opt{A}{(c) and (d) only}
\opt{B}{(a), (c) and (d) only}
\opt{C}{(b), (c) and (d) only}
\opt{D}{(a), (b), (c) and (d)}
\end{qoptions}
\qn{19} Which of the following statements are correct ?
\par\noindent(a) Both LiCl and MgCl$_2$ are soluble in ethanol.\par
\noindent(b) The oxides Li$_2$O and MgO combine with excess of oxygen to give superoxide.\par
\noindent(c) LiF is less soluble in water than other alkali metal fluorides.\par
\noindent(d) Li$_2$O is more soluble in water than other alkali metal oxides.\par
Choose the most appropriate answer from the options given below:
\begin{qoptions}
\opt{A}{(a) and (c) only}
\opt{B}{(a), (c) and (d) only}
\opt{C}{(b) and (c) only}
\opt{D}{(a) and (c) only}
\end{qoptions}
\qn{20} The initial rate of hydrolysis of methyl acetate (1M) by a weak acid (HA, 1M) is $1/100^{\text{th}}$ of that of a strong acid (HX, 1M), at 25$^\circ$C. The $K_a$ of HA is
\begin{qoptions}
\opt{A}{$1\times10^{-4}$}
\opt{B}{$1\times10^{-5}$}
\opt{C}{$1\times10^{-6}$}
\opt{D}{$1\times10^{-3}$}
\end{qoptions}
\sect{SECTION-II}
\qn{1} How many reaction correctly match with their major product :
\qfig{p2_chem_s2_q01_1}
\qn{2} \qfig{p2_chem_s2_q02_1}
Number of CO$_2$ evolved.
\qn{3} 50 mL of 0.1 M CH$_3$COOH is being titrated against 0.1 M NaOH. When 25 mL of NaOH has been added, the pH of the solution will be \underline{\hspace{1.2cm}} $\times10^{-2}$. (Nearest integer)
\par\noindent(Given : p$K_a$ (CH$_3$COOH) = 4.76)\par
\noindent$\log2=0.30$\par\noindent$\log3=0.48$\par\noindent$\log5=0.69$\par\noindent$\log7=0.84$\par\noindent$\log11=1.04$
\qn{4} Find $\Delta G$ (J/mol) for the reaction at 300 kPa \& $27^\circ$C when all gases are in stoichiometric ratio of moles. N$_2$O$_4$(g) $\to$ 2NO$_2$(g) [$\ln2=0.7$]
\par\noindent Given : $\Delta G^0_f(\mathrm{N_2O_4})=100$ kJ/mol\par
\noindent\hspace{3.2em}$\Delta G^0_f(\mathrm{NO_2})=50$ kJ/mol\par
\noindent\hspace{3.2em}$R=8$ J/mol-K
\qn{5} A solution containing $2.5\times10^{-3}$ kg of a solute dissolved in $75\times10^{-3}$ kg of water boils at 373.535 K. The molar mass of the solute is \underline{\hspace{1.2cm}} g mol$^{-1}$. [nearest integer] (Given: $K_b$(H$_2$O) = 0.52 K Kg mol$^{-1}$, boiling point of water = 373.15K)
\qn{6} Identify total number of compounds which give both tollen's and fehling test ?
\par\noindent(i) HCHO\quad(ii) CH$_3$--CHO\quad(iii) Ph--CHO\quad(iv) CH$_3$--C$\equiv$C--CH$_3$\par
\noindent(v) \qfigi{p2_chem_s2_q06_2}\quad(vi) Ph--CH$_2$--CHO\par
\noindent(vii) HCOOH\quad(viii) CH$_3$--CO--CH$_2$--OH\par
\noindent(ix) CH$_3$--COOH\quad(x) Ph--CO--Ph
\qn{7} Match given properties with elements
\begin{center}\footnotesize\setlength{\tabcolsep}{2pt}
\begin{tabular}{|L{0.35cm}L{2.9cm}|L{0.35cm}L{1.4cm}|}\hline
\multicolumn{2}{|c|}{List-I (Properties)}&\multicolumn{2}{c|}{List-II (Elements)}\\\hline
P&Produce stable M$_2$O type oxide with limited supply of air&1&Li, Be, Al\\\hline
Q&Chloride of elements produce red colour in oxidising flame&2&Li, Ca, Sr\\\hline
R&Chlorides of such elements having significant covalent character or covalent nature in anhydrous form&3&Na, Li, Cs\\\hline
\end{tabular}\end{center}
Match list-I and list-II and give your answer as abc where a,b and c are the match of 1, 2, 3 from list-II.

\textbf{(for example Matching of P is = 2, Q = 3 \& R = 1 then your answer is 231)}
\qn{8} \qfig{p2_chem_s2_q08_1}
\par\noindent(i) FeSO$_4$ + KCN(excess) $\to$ (X) \par\noindent\hspace{1.5em}(X) $\xrightarrow{\text{conc. H}_2\text{SO}_4}$ (A) gas\hspace{1em}(X) $\xrightarrow{\text{dil. H}_2\text{SO}_4}$ (B) gas
\par\noindent(ii) CaC$_2$ + N$_2$ $\xrightarrow{1100^\circ\text{C}}$ (Y) + Carbon
\par Find the sum of the number of $\pi$ bonds present in (A), (B) and (Y).
\qn{9} Match the colour with given compounds
\begin{center}\footnotesize\setlength{\tabcolsep}{2pt}
\begin{tabular}{|L{0.35cm}L{1.8cm}|L{0.35cm}L{1.8cm}|}\hline
\multicolumn{2}{|c|}{List-I (Compounds)}&\multicolumn{2}{c|}{List-II (Colour)}\\\hline
P&Mn$_2$O$_7$&1&Orange\\\hline
Q&(NH$_4$)$_2$Cr$_2$O$_7$&2&Green\\\hline
R&CrO$_2$Cl$_2$&3&Reddish Brown\\\hline
\end{tabular}\end{center}
Match list-I and list-II and give your answer as abc where a,b and c are the match of 1, 2, 3 from list-II.

\textbf{(for example Matching of P is = 2, Q = 3 \& R = 1 then your answer is 231)}
\qn{10} \begin{center}\footnotesize\setlength{\tabcolsep}{2pt}
\begin{tabular}{|L{0.6cm}L{1.5cm}|L{0.6cm}L{2.2cm}|}\hline
\multicolumn{2}{|c|}{List-I Catalyst}&\multicolumn{2}{c|}{List-II Process}\\\hline
(P)&V$_2$O$_5$&(1)&Ostwald process\\\hline
(Q)&CuCl$_2$&(2)&Deacon process\\\hline
(R)&Pt/Rh&(3)&Contact process\\\hline
\end{tabular}\end{center}
Match list-I and list-II and give your answer as abc where a,b and c are the match of 1, 2, 3 from list-II.

\textbf{(for example Matching of P is = 2, Q = 3 \& R = 1 then your answer is 231)}
\columnbreak
\partsec{3}{MATHEMATICS}
\sect{SECTION-I}
\qn{1} If the function $f$, $g$, $h$ are defined from the sets of real numbers $\mathbb R$ to $\mathbb R$ such that $f(x)=x^2-1$, $g(x)=\sqrt{x^2+1}$, $h(x)=\begin{cases}0,&\text{if } x\le0\\ x,&\text{if } x>0\end{cases}$, then the composite function $h\circ f\circ g(x)$ is
\begin{qoptions}
\opt{A}{$\begin{cases}0,&x=0\\ x^2,&x>0\\ -x^2,&x<0\end{cases}$}
\opt{B}{$\begin{cases}0,&x=0\\ x^2,&x\ne0\end{cases}$}
\opt{C}{$\begin{cases}0,&x<0\\ x^2,&x\ge0\end{cases}$}
\opt{D}{$\begin{cases}x^2,&x\le0\\ -x^2,&x>0\end{cases}$}
\end{qoptions}
\qn{2} If $\vec a,\vec b$ are non-zero vectors such that $\left|\vec a+\vec b\right|=\left|\vec a-2\vec b\right|$, then -
\begin{qoptions}
\opt{A}{$\vec a\cdot\vec b=2\left|\vec b\right|^2$}
\opt{B}{$\vec a\cdot\vec b=\left|\vec b\right|^2$}
\opt{C}{Least value of $\vec a\cdot\vec b+\dfrac{4}{\left|\vec b\right|^2+2}$ is $2\sqrt2$}
\opt{D}{Least value of $\vec a\cdot\vec b+\dfrac{4}{\left|\vec b\right|^2+2}$ is $2\sqrt2-1$}
\end{qoptions}
\qn{3} If $A=\begin{bmatrix}1&0&0\\0&1&1\\0&-2&4\end{bmatrix}$ ; $6A^{-1}=A^2+cA+dI$ then $(c,d)$ is
\begin{qoptions}
\opt{A}{$(-6,11)$}
\opt{B}{$(-11,6)$}
\opt{C}{$(11,6)$}
\opt{D}{$(6,11)$}
\end{qoptions}
\qn{4} Unit vectors $\vec a,\vec b$ and $\vec c$ are coplanar, unit vector $\vec d$ is perpendicular to them. If $\left(\vec a\times\vec b\right)\times\left(\vec c\times\vec d\right)=\dfrac{\hat i-2\hat j+2\hat k}{6}$ and angle between $\vec a$ and $\vec b$ is $30^\circ$ then $\vec c$ can be
\begin{qoptions}
\opt{A}{$\dfrac{\hat i-\hat j+\hat k}{\sqrt3}$}
\opt{B}{$\dfrac{\hat i-2\hat j+2\hat k}{3}$}
\opt{C}{$\dfrac{2\hat i-\hat j+\hat k}{\sqrt6}$}
\opt{D}{$\dfrac{\hat i-2\hat j+2\hat k}{6}$}
\end{qoptions}
\qn{5} Consider the three linear equations with respect to $X$, $Y$ and $Z$ as
\par\noindent$x.X+Y+2Z=0$\par\noindent$f(x).X+3Y+x^2Z=0$\par\noindent$5xX+6y+Z=0$ having
\par\noindent non-trivial solution for $X$, $Y$ and $Z$ then curve $f(x)$
\begin{qoptions}
\opt{A}{is always increasing}
\opt{B}{is always decreasing}
\opt{C}{has exactly one critical point}
\opt{D}{has three critical points}
\end{qoptions}
\qn{6} The value of definite integral
\par\noindent\hfill$\displaystyle\begin{aligned}\int_{-\pi}^{\pi}\bigl(&\cos2x\cdot\cos2^2x\cdot\cos2^3x\\&\cdot\cos2^4x\cdot\cos2^5x\bigr)dx\end{aligned}$\hfill\null\par
is
\begin{qoptions}
\opt{A}{1}
\opt{B}{$-1$}
\opt{C}{2}
\opt{D}{0}
\end{qoptions}
\qn{7} If $\theta_1=\sin^{-1}\!\left(\dfrac45\right)+\sin^{-1}\!\left(\dfrac13\right)$ and $\theta_2=\cos^{-1}\dfrac45+\cos^{-1}\dfrac13$ then
\begin{qoptions}
\opt{A}{$\theta_1<\theta_2$}
\opt{B}{$\theta_1>\theta_2$}
\opt{C}{$\theta_1=\theta_2$}
\opt{D}{$\theta_1-2\theta_2$}
\end{qoptions}
\qn{8} The function $f:[a,\infty)\to R$, where $R$ denotes the range of $f(x)$ corresponding to the given domain $[a,\infty)$. If $f(x)=2x^3-3x^2+6$ will have an inverse, then complete set of possible values of $a$ is
\begin{qoptions}
\opt{A}{$(-\infty,1]$}
\opt{B}{$[0,\infty)$}
\opt{C}{$(-\infty,0]$}
\opt{D}{$[1,\infty)$}
\end{qoptions}
\qn{9} If $f(x)=\displaystyle\int e^x\left(\tan^{-1}(x)+\dfrac{2x}{(1+x^2)^2}\right)dx$, $f(0)=-1$ then the value of $f(1)$ is
\begin{qoptions}
\opt{A}{$e\left(\dfrac\pi4-\dfrac12\right)$}
\opt{B}{$e\left(\dfrac\pi4+\dfrac12\right)$}
\opt{C}{$e\left(\dfrac\pi2-\dfrac14\right)$}
\opt{D}{$e^{-1}\left(\dfrac\pi4-\dfrac12\right)$}
\end{qoptions}
\qn{10} If $\sqrt{x+y}+\sqrt{y-x}=c$ then $\dfrac{d^2y}{dx^2}$ equals
\begin{qoptions}
\opt{A}{$\dfrac2{c^2}$}
\opt{B}{$\dfrac{-2}{c^2}$}
\opt{C}{$\dfrac2c$}
\opt{D}{$\dfrac{-2}c$}
\end{qoptions}
\qn{11} If $I_n=\displaystyle\int(\operatorname{cosec}2x)^n\,dx$ then
\begin{qoptions}
\opt{A}{$I_n=\dfrac{\cot2x\operatorname{cosec}^{n-2}2x}{-2(n-1)}+\dfrac{n-2}{n-1}I_{n-2}$}
\opt{B}{$I_n=\dfrac{\cot2x\operatorname{cosec}^{n-2}2x}{-2(n-1)}+\dfrac{n+2}{n-1}I_{n-2}$}
\opt{C}{$I_n=\dfrac{\cot2x\operatorname{cosec}^{n-2}2x}{-(n-1)}+\dfrac{n-2}{n-1}I_{n-2}$}
\opt{D}{$I_n=\dfrac{\cot2x\operatorname{cosec}^{n-2}2x}{-(n-1)}+\dfrac{n+2}{n-1}I_{n-2}$}
\end{qoptions}
\qn{12} If water is poured into an inverted hollow right circular cone whose semi-vertical angle is $30^\circ$, its depth increase at the rate of 1 cm/s. The rate at which the volume of water increase when the depth is 24 cm is equal to :
\begin{qoptions}
\opt{A}{$150\pi$ cm$^3$/s}
\opt{B}{$180\pi$ cm$^3$/s}
\opt{C}{$192\pi$ cm$^3$/s}
\opt{D}{$200\pi$ cm$^3$/s}
\end{qoptions}
\qn{13} The minimum distance between the curves $y=\tan x$, $x\in\left(-\dfrac\pi2,\dfrac\pi2\right)$ and $\left(x-2-\dfrac\pi4\right)^2+y^2=1$ is
\begin{qoptions}
\opt{A}{$\sqrt5+1$}
\opt{B}{$\sqrt5-1$}
\opt{C}{$\sqrt2-1$}
\opt{D}{2}
\end{qoptions}
\qn{14} If $f(x)=1-x-x^3$, then the set of all real values of $x$ satisfying the inequality $1-f(x)-f^3(x)>f(1-5x)$ is
\begin{qoptions}
\opt{A}{$(-2,0)\cup(2,\infty)$}
\opt{B}{$(-\infty,-2)$}
\opt{C}{$(1,2)$}
\opt{D}{$(0,2)$}
\end{qoptions}
\qn{15} Area bounded by the curves $y=-x^2+6x-5$, $y=-x^2+4x-3$ and $y=3x-15$ is
\begin{qoptions}
\opt{A}{$\dfrac{73}2$}
\opt{B}{$\dfrac{73}3$}
\opt{C}{$\dfrac{73}6$}
\opt{D}{73}
\end{qoptions}
\qn{16} If $\displaystyle\sum_{r=0}^{n}\dfrac{(-1)^r\cdot{}^nC_r}{(r+1)(r+2)(r+3)}=\dfrac1{a(n+b)}$, then $a+b$ is equal to
\begin{qoptions}
\opt{A}{2}
\opt{B}{4}
\opt{C}{5}
\opt{D}{6}
\end{qoptions}
\qn{17} If $\dfrac2{1!9!}+\dfrac2{3!7!}+\dfrac1{5!5!}=\dfrac{2^m}{n!}$, then which of the following is correct ? (Where $m,n\in N$)
\begin{qoptions}
\opt{A}{$n-m=1$}
\opt{B}{$n+m=18$}
\opt{C}{$m-n=1$}
\opt{D}{$m+n=20$}
\end{qoptions}
\qn{18} $\displaystyle\lim_{x\to0^+}\left((\sin x)^{\frac1x}+(\cos x)^{\frac1x}+(\operatorname{cosec}x)^x\right)$ is equal to
\begin{qoptions}
\opt{A}{2}
\opt{B}{4}
\opt{C}{6}
\opt{D}{8}
\end{qoptions}
\qn{19} If $y(x)$ is the solution of the differential equation $\dfrac{dy}{dx}=-2x(y-1)$ with $y(0)=1$, then $\displaystyle\lim_{x\to\infty}y(x)$ equals
\begin{qoptions}
\opt{A}{2}
\opt{B}{1}
\opt{C}{4}
\opt{D}{5}
\end{qoptions}
\qn{20} Let $\vec v$ be a vector which satisfies the condition $2\vec v+\vec v\times\left(\hat i+2\hat j\right)=2\hat i+\hat k$. If $3\left|\vec v\right|=\sqrt m$, then find the value of $m$ ?
\begin{qoptions}
\opt{A}{8}
\opt{B}{10}
\opt{C}{6}
\opt{D}{7}
\end{qoptions}
\sect{SECTION-II}
\qn{1} Let $m$ be the minimum value of the function $f(x)=2x^3-3x^2-12x+6$ on the set $A=\{x:\ell nx>\ell n(x^2-3x+3)\}$, then find $|m|$
\qn{2} Let $f:R\to\left(0,\dfrac\pi2\right]$ defined by $f(x)=\operatorname{cosec}^{-1}(ax^2-2px+4)$ where $a,p\in N$ is surjective then smallest value of $2(a+p)$ is
\qn{3} Let $A$ is a non-zero matrix of order `2' such that $A+\operatorname{adj}A=A^{-1}$ then determinant of $(2A^{-1})$ is equal to
\qn{4} Let $[a_{ij}]$ is a square matrix of order 2 where $a_{ij}\in\{0,1,2,3,4,6\}$ then number of matrices $A$ such that $A^{-1}$ exists is $\lambda$ then $\lambda-300$ is (given that all elements of matrix $A$ are distinct)
\qn{5} Let $\vec a=2\hat i+\hat k$, $\vec b=\hat i+\hat j+\hat k$ and $\vec c=4\hat i+3\hat j+7\hat k$. If $\vec r\times\vec b=\vec c\times\vec b$ and $\vec r\cdot\vec b=\vec c\cdot\vec b$, then the value of $\vec a\cdot\vec r=$
\qn{6} If $|\vec a|=|\vec b|=|\vec c|=1$ and $\vec a\mathbin{\wedge}\vec b=\vec c\mathbin{\wedge}\vec b=\vec a\mathbin{\wedge}\vec c=\theta$ where $\vec p\mathbin{\wedge}\vec q$ represents the angle between vectors $\vec p$ and $\vec q$ then maximum value of $\theta$ is $\dfrac\pi k$. The value of $4k^2$ is
\qn{7} $I=\displaystyle\int_0^2(3x^2-3x+1)\cos\left(x^3-3x^2+4x-2\right)dx=p(\sin q)$ where $p,q\in N$, then find $(p^2+q^3)$.
\qn{8} $\displaystyle\sum_{i=0}^{100}{}^{150}C_i\cdot{}^{350}C_{100-i}\,(150-i)=k\,{}^{500}C_{100}$ then value of $k$ is equal to
\qn{9} The last digit of $3^{81^n}+1$, $n\in N$, is
\qn{10} If $\vec a$ and $\vec b$ are non collinear vectors such that $(2x-y)\vec a+4\vec b=5\vec a+(x-2y)\vec b$ then $(x+y)$ is
\end{multicols}
\end{document}
